% from ega4-18-fr.tex
\oldpage[IV-4]{151}

dans $X$), $g^{-1}(U)$ est connexe (Bourbaki, \emph{Top. gén.}, chap.~I\textsuperscript{er}, 3\textsuperscript{e} éd., \textsection~11, n\textsuperscript{o}~3, prop.~7). Si l'on désigne par $\pi_0(X)$ l'ensemble des composantes connexes de $X$, l'application $\pi_0(X')\to\pi_0(X)$ déduite canoniquement de $g$ est donc \emph{bijective}.

\begin{corollary}[18.9.5]
\label{IV.18.9.5}
Sous les hypothèses de (18.9.4), le foncteur
\begin{equation}
  Z\longmapsto Z\otimes_A A'
  \tag{18.9.5.1}
  \label{IV.18.9.5.1}
\end{equation}
de la catégorie des préschémas étales sur $X$, dans celle des préschémas étales sur $X'$, est pleinement fidèle.
\end{corollary}

Soient $Z_1$, $Z_2$ deux préschémas étales sur $X$, et posons $Z'_i=Z_i\otimes_A A'$ ($i=1,2$); il s'agit de prouver que tout $X'$-morphisme $Z'_1\to Z'_2$ provient, par changement de base, d'un $X$-morphisme unique $Z_1\to Z_2$. Supposons d'abord que $Z_2$ soit \emph{séparé} sur $X$; alors, comme $\Hom_X(Z_1,Z_2)$ s'identifie à $\Gamma((Z_1\times_X Z_2)/Z_1)$, et que $Z_1\times_X Z_2$ est étale et séparé sur $Z_1$, $\Hom_X(Z_1,Z_2)$ s'identifie fonctoriellement, en vertu de (17.9.3), à l'ensemble des parties ouvertes et fermées $U$ de $Z_1\times_X Z_2$ telles que la restriction à $U$ de la projection $p:Z_1\times_X Z_2\to Z_1$ soit un morphisme surjectif et radiciel. L'assertion résulte donc dans le cas considéré de (18.9.4) et de ce que $Z'_1\times_{X'}Z'_2=(Z_1\times_X Z_2)\otimes_A A'$.

Passons au cas général: il suffira de prouver que le graphe $\Gamma'$ d'un $X'$-morphisme $Z'_1\to Z'_2$, qui est ouvert (17.9.3) dans $Z'_1\times_{X'}Z'_2$, est de la forme $\Gamma\otimes_A A'$, où $\Gamma$ est induit sur une partie \emph{ouverte} de $Z=Z_1\times_X Z_2$. En effet, la restriction à $\Gamma$ de la projection $p:Z\to Z_1$ est alors un \emph{isomorphisme}, car par changement de base $A\to A'$ cette restriction devient l'isomorphisme restriction à $\Gamma'$ de la projection $p':Z'\to Z'_1$ (avec $Z'=Z\otimes_A A'$), et il suffit d'appliquer (2.7.1, (viii)); la conclusion résulte alors de la caractérisation des graphes de morphismes (I, 5.3).

Si $q:Z'\to Z$ est la projection canonique, pour prouver que $\Gamma'$ est de la forme $q^{-1}(\Gamma)$, où $\Gamma$ est ouvert dans $Z$, il suffit, puisque $q$ est un morphisme fidèlement plat et quasi-compact, de prouver qu'il existe un ensemble $U\subset Z$ tel que $\Gamma'=q^{-1}(U)$ (2.3.12). Posons $Z''=Z'\times_Z Z'=Z\times_X X''$, où $X''=X'\times_X X'$, et soient $q_1:Z''\to Z'$, $q_2:Z''\to Z'$ les projections canoniques. Appliquant (4.5.19.1), il suffit de montrer que $q_1^{-1}(\Gamma')=q_2^{-1}(\Gamma')$. Mais c'est une propriété qui est vraie si et seulement si elle l'est après chaque changement de base $\Spec(k(x))\to X$, où $x$ parcourt $X$. Autrement dit, il suffit de prouver le corollaire lorsque $X$ est le \emph{spectre d'un corps}; mais comme tout $X$-préschéma étale est alors automatiquement \emph{séparé} sur $X$ (17.6.2, $c'$)), on est ramené au cas envisagé au début de la démonstration.

\begin{remarks}[18.9.6]
\label{IV.18.9.6}
\begin{enumerate}
  \item[(i)] Il est possible que, si le corps résiduel de $A$ est de caractéristique 0, le foncteur (18.9.5.1) induise même une \emph{équivalence} de la catégorie des revêtements étales de $X$ et de la catégorie des revêtements étales de $X'$. On peut en tout cas le démontrer lorsque $A$ est \emph{excellent}, en utilisant la résolution des singularités de Hironaka (M.~Artin). Il est plausible que l'énoncé analogue soit encore vrai sans restriction sur la caractéristique du corps résiduel, à condition de se borner aux revêtements « galoisiens principaux » dont le groupe est d'ordre premier à la caractéristique résiduelle (cf. [41]).

